\magicSection{Tensors}{sec:Tensors}

I kind of hate the debasement of the mathematical term ``tensor''.
Anyway a tensor has
\begin{itemize}
  \item An element type (float32, float16, etc.)
  \item Physical memory type (registers, shared memory, global memory, CPU memory)
  \item Allocator (e.g. stack vs malloc allocation. Merged into memory type in Exo, but fundamentally separate).
    This may impose synchronization requirements upon alloc and free, in particular for shared memory (not further discussed here).
  \item Per-dimension \myKeyA{extent} and \myKeyA{stride}
  \item Opaque swizzle function for shared memory allocations (Section~\ref{sec:StridingAndSwizzling})
\end{itemize}
The strides are the interesting part.
Each stride is a multiple of a certain ``stride base unit''
\begin{itemize}
  \item Bits, for packing multiple elements into words
  \item RAM address, for tensors not stored in registers
  \item Register ID, for tensors stored in registers
  \item Thread ID, for tensors sharded across threads
  \item CTA ID, for tensors in distributed shared memory
\end{itemize}
Let's disect the format of the row-major A matrix for float16 mma instructions,
using temporary syntax \lighttt{extent @ stride} for each dimension.

The base case is a single 32-bit word packing 2 16-bit floats

\input{ultraspork/mma_a.0.tex}

A single row contains 8 16-bit floats.
This is sharded across 4 threads, with 2 packed values per thread.

\input{ultraspork/mma_a.1.tex}

From there, a whole warp (32 threads) stores a single $8 \times 8$ matrix.
Each group of 4 threads stores one row.

\input{ultraspork/mma_a.2.tex}

Thus far only a single register ID was consumed.
The full instruction takes a $16 \times 16$ $A$ matrix, split into 4 registers.
Each register stores a single $8 \times 8$ matrix, with adjacent matrices on the $M$ dimension having adjacent registers.

\input{ultraspork/mma_a.3.tex}

i.e., the stride on the $M$ register dimension is 1 register ID, and on the $K$ register dimension, 2 register IDs.

This implies a polynomial indexing function mapping tensor coordinates to an ``address'' (I need an abstract term for a hardware location that is not necessarily a pointer).
There are free variables for the unknown base of the allocation, and the stride base units.

\input{ultraspork/mma_a.4.tex}

For clarity we should consider CuTe-style hierarchical coordinates (tuple of tuples).
In this example, group the $M$ and $K$ coordinates separately.
We can use lexicographical ordering to implement what CuTe calls ``natural coordinates'' where a single coordinate is auto-decomposed to a coordinate tuple.
The implied indexing function is now polynomial, with the addition of an opaque function \texttt{DivAndMod(x, D, M) = $\frac{x}{D} \% M$}.

\input{ultraspork/mma_a.5.tex}

